% =========================================================================
% SciPost LaTeX template — Lecture notes
% =========================================================================

\documentclass{SciPost}

% Prevent all line breaks in inline equations.
\binoppenalty=10000
\relpenalty=10000

\hypersetup{
    colorlinks,
    linkcolor={red!50!black},
    citecolor={blue!50!black},
    urlcolor={blue!80!black}
}

\usepackage[bitstream-charter]{mathdesign}
\urlstyle{same}

% Number lines in verbatim code blocks for easy reference
\usepackage{fancyvrb}
\RecustomVerbatimEnvironment{verbatim}{Verbatim}{numbers=left,numbersep=5pt,frame=single,fontsize=\small}

% Fix \cal and \mathcal characters look (so it's not the same as \mathscr)
\DeclareSymbolFont{usualmathcal}{OMS}{cmsy}{m}{n}
\DeclareSymbolFontAlphabet{\mathcal}{usualmathcal}

\fancypagestyle{SPstyle}{
\fancyhf{}
\lhead{\colorbox{scipostblue}{\bf \color{white} ~SciPost Physics Lecture Notes }}
\rhead{{\bf \color{scipostdeepblue} ~Draft }}
\renewcommand{\headrulewidth}{1pt}
\fancyfoot[C]{\textbf{\thepage}}
}

% Additional packages not provided by SciPost
\usepackage{mathtools}
\usepackage{physics}
\usepackage{cleveref}
\usepackage{enumitem}
\usepackage{booktabs}
\usepackage{array}
\usepackage{tikz-cd}
\usepackage{amsthm}

% Theorem environments
\newtheorem{theorem}{Theorem}[section]
\newtheorem{proposition}[theorem]{Proposition}
\newtheorem{lemma}[theorem]{Lemma}
\newtheorem{corollary}[theorem]{Corollary}
\theoremstyle{definition}
\newtheorem{definition}[theorem]{Definition}
\newtheorem{example}[theorem]{Example}
\newtheorem{remark}[theorem]{Remark}

% Shortcuts
\newcommand{\Z}{\mathbb{Z}}
\newcommand{\Zt}{\mathbb{Z}_2}
\newcommand{\Zn}[1]{\mathbb{Z}_{#1}}
\newcommand{\U}{\mathrm{U}}
\newcommand{\Gb}{G_b}
\newcommand{\Gf}{G_f}
\newcommand{\dd}{\mathrm{d}}
\renewcommand{\im}{\operatorname{im}}
\newcommand{\coker}{\operatorname{coker}}
\newcommand{\gr}{\operatorname{gr}}
\renewcommand{\Hom}{\operatorname{Hom}}
\newcommand{\Ext}{\operatorname{Ext}}


\begin{document}

\pagestyle{SPstyle}

\begin{center}{\Large \textbf{\color{scipostdeepblue}{
Landscape of Interacting Fermionic Symmetry-Protected Topological Phases
}}}\end{center}

\begin{center}\textbf{
Xing-Yu Ren\textsuperscript{1$\star$}
}\end{center}

\begin{center}
{\bf 1} Department of Physics, The Chinese University of Hong Kong
\\ [\baselineskip]
$\star$ \href{mailto:xy.ren@link.cuhk.edu.hk}{\small xy.ren@link.cuhk.edu.hk}\
\end{center}


\section*{\color{scipostdeepblue}{Abstract}}
\textbf{\boldmath{%
Symmetry-protected topological (SPT) phases reveal deep connections between symmetry and quantum entanglement.
While free-fermion classifications are well established, interacting fermionic SPT phases---especially under crystalline symmetries---exhibit a far richer structure.
In this lecture, I present a systematic landscape based on domain wall decoration and organized by the Atiyah--Hirzebruch spectral sequence (AHSS), yielding a complete classification of fermionic SPT phases protected by space group symmetries and uncovering new interacting invariants beyond free-fermion approaches.
These lecture notes are designed to be self-contained for physicists who may not have encountered spectral sequences before, with emphasis on the physical intuition behind the mathematical formalism.
}}

\linenumbers

\vspace{10pt}
\noindent\rule{\textwidth}{1pt}
\tableofcontents
\noindent\rule{\textwidth}{1pt}
\vspace{10pt}

%% ============================================================
%% CHAPTER 1
%% ============================================================
\section{Introduction}
\label{sec:intro}

\subsection{What is a topological phase?}

Consider a quantum many-body system at zero temperature, described by a gapped Hamiltonian $H$.
Two gapped Hamiltonians $H_0$ and $H_1$ are in the \emph{same phase} if there exists a smooth path of gapped Hamiltonians $H(t)$, $t\in[0,1]$, connecting them without closing the energy gap.
A \emph{topological phase} is an equivalence class under this relation.

The trivial phase is the one containing a product state (a state with no entanglement at all).
All other gapped phases are ``topological'' in the sense that they carry patterns of quantum entanglement that cannot be smoothly removed.

\subsection{Symmetry protection}

Many interesting topological phases become trivial if we do not impose symmetry.
A \emph{symmetry-protected topological} (SPT) phase is a gapped phase that:
\begin{enumerate}[label=(\roman*)]
    \item has a unique ground state on any closed manifold (no topological order),
    \item is distinct from the trivial phase \emph{only when} a specified symmetry $G$ is preserved along the interpolation path.
\end{enumerate}
The most familiar examples include the Haldane chain (protected by $\Zt \times \Zt$ or $SO(3)$) and topological insulators (protected by time-reversal and charge conservation).

\subsection{Free fermions vs.\ interactions}

For non-interacting (free) fermion systems, the classification of topological phases is completely understood through $K$-theory \cite{Kitaev2009, Ryu2010, Schnyder2008}.
The result is the celebrated \emph{periodic table} of topological insulators and superconductors, organized by spatial dimension and the ten Altland--Zirnbauer symmetry classes.

However, this classification assumes the absence of interactions.
When interactions are included, two qualitatively new phenomena occur:
\begin{enumerate}[label=(\roman*)]
    \item \textbf{Interaction-induced collapse}: Some phases that are distinct in the free-fermion classification become equivalent once interactions are allowed.
    A celebrated example is the $\Z$ classification of 1D class BDI topological superconductors, which collapses to $\Z_8$ upon including interactions \cite{Fidkowski2010, Fidkowski2011}.
    \item \textbf{Interaction-enabled phases}: New phases can appear that have no free-fermion counterpart.
    These are ``intrinsically interacting'' topological phases.
\end{enumerate}
The goal of this lecture is to present a systematic framework that captures \emph{all} interacting fermionic SPT phases for a given symmetry, not just those visible to band theory.

\subsection{Our approach: decoration + spectral sequence}

The framework we use has two complementary aspects:
\begin{enumerate}[label=(\roman*)]
    \item \textbf{Domain wall decoration} (physics): A constructive method to build SPT states by decorating symmetry domain walls with lower-dimensional topological phases.
    This gives a concrete physical picture for every phase in the classification.
    \item \textbf{Atiyah--Hirzebruch spectral sequence} (mathematics): A powerful algebraic machine that organizes the decoration data, systematically accounts for obstructions and equivalences, and yields the classification as a computable algebraic object.
\end{enumerate}
The deep insight is that these two aspects are \emph{the same thing viewed from different angles}: the physical construction naturally produces a filtered cochain complex, and AHSS is precisely the standard mathematical tool for computing the cohomology of filtered complexes.

Moreover, the entire framework has been implemented as a computational tool---the \texttt{SptSet} package for the GAP computer algebra system---allowing systematic, automated classification for any symmetry group.

\subsection{Plan of the lecture}

The lecture is organized into two main parts:

\medskip
\noindent\textbf{Part I: Classification.}
\begin{itemize}
    \item \S\ref{sec:decoration}: The physics of domain wall decoration---how to construct FSPT states layer by layer, and the consistency conditions (coboundary equations) that the decoration data must satisfy.
    \item \S\ref{sec:AHSS}: The AHSS formalism---how the layered decoration structure maps into a spectral sequence, and how the spectral sequence systematically computes the classification. \emph{This is the mathematical heart of the lecture.}
    \item \S\ref{sec:sptset-classification}: Hands-on use of the \texttt{SptSet} package to compute classifications.
\end{itemize}

\medskip
\noindent\textbf{Part II: Stacking.}
\begin{itemize}
    \item \S\ref{sec:stacking-physics}: The physics of stacking---what happens when you combine two SPT phases, and why the group law has nontrivial cross-terms.
    \item \S\ref{sec:extension}: The extension problem in AHSS---how stacking maps to the mathematical problem of reconstructing a group from its graded pieces.
    \item \S\ref{sec:sptset-stacking}: Computing stacking groups with \texttt{SptSet}.
\end{itemize}

\medskip
\noindent We conclude in \S\ref{sec:outlook} with results, comparisons to free-fermion classifications, and open problems.


%% ============================================================
%% CHAPTER 2
%% ============================================================
\section{Domain Wall Decoration: Physics of Classification}
\label{sec:decoration}

This section introduces the physical framework of domain wall decoration for constructing fermionic SPT phases.
The key idea is simple: decorate symmetry domain walls with lower-dimensional topological phases, then impose consistency conditions.

\subsection{Fermionic symmetry}
\label{subsec:fermionic-symmetry}

In a fermionic system, the fermion parity symmetry $\Zt^f$ is always present and cannot be explicitly broken.
The total symmetry group $\Gf$ is a central extension of the bosonic symmetry group $\Gb$ by the fermion parity:
\begin{equation}
\label{eq:Gf-extension}
1 \longrightarrow \Zt^f \longrightarrow \Gf \longrightarrow \Gb \longrightarrow 1.
\end{equation}
This extension is classified by a 2-cocycle $\omega_2 \in H^2(\Gb, \Zt)$.
Physically, $\omega_2$ encodes whether the symmetry operators of $\Gb$ commute or anticommute with the fermion parity.

If the system has anti-unitary symmetries (such as time-reversal), we additionally specify a group homomorphism
\begin{equation}
\label{eq:s1}
s_1 \in \Hom(\Gb, \Zt) = H^1(\Gb, \Zt),
\end{equation}
where $s_1(g) = 0$ if $g$ is unitary and $s_1(g) = 1$ if $g$ is anti-unitary.

\begin{definition}[Fermionic symmetry]
The symmetry of an internal fermionic SPT phase is specified by a triplet
\begin{equation}
\Gf := (\Gb,\, \omega_2,\, s_1).
\end{equation}
\end{definition}

\begin{example}
Consider the symmetry group $\Gf = \Zn{4}^{f,T}$, generated by time-reversal $T$ with $T^2 = (-1)^F$ (the fermion parity).
Then $\Gb = \Zt^T$ with $s_1(T) = 1$ (anti-unitary) and $\omega_2(T, T) = 1$ (nontrivial extension).
\end{example}


\subsection{The decoration scheme}
\label{subsec:decoration-scheme}

The construction of fermionic SPT states proceeds on a triangulated spatial manifold equipped with a branching structure (a consistent choice of edge orientations that avoids oriented loops in each simplex).

The key physical idea is to break the symmetry $\Gb$ explicitly, creating domain walls between different symmetry-breaking regions, and then decorate these domain walls with lower-dimensional invertible topological phases.
Afterwards, one restores the symmetry by quantum superposing all possible symmetry-breaking patterns and their decorations.

For a $d$-dimensional spatial system, the decoration data consists of \emph{layers} of topological phases living on simplices of different dimensions:

\begin{center}
\renewcommand{\arraystretch}{1.3}
\begin{tabular}{cccl}
\toprule
\textbf{Layer} & \textbf{Cochain} & \textbf{Coefficient} & \textbf{Physical content} \\
\midrule
$d$ (top) & $\nu_{d+1}$ & $\U(1)$ & Bosonic SPT phase factor \\
$d-1$ & $n_d$ & $\Zt$ & Complex fermion decoration \\
$d-2$ & $n_{d-1}$ & $\Zt$ & Majorana chain decoration \\
$d-3$ & $n_{d-2}$ & $\Z$ & $p+ip$ superconductor decoration \\
\bottomrule
\end{tabular}
\end{center}

Each layer is described by an \emph{inhomogeneous cochain} on $\Gb$ with values in the corresponding coefficient module.
Specifically, a $p$-cochain $f \in C^p(\Gb, A)$ is a function
\begin{equation}
f: \underbrace{\Gb \times \cdots \times \Gb}_{p} \longrightarrow A, \qquad (g_1, \ldots, g_p) \mapsto f(g_1, \ldots, g_p),
\end{equation}
that assigns a value in the coefficient group $A$ to each $p$-tuple of group elements.

\subsection{Consistency conditions: the coboundary equations}
\label{subsec:coboundary-eqs}

The decoration data cannot be chosen freely.
When the symmetry-breaking pattern fluctuates (as required by the restoration of symmetry), the decorations on adjacent simplices must be \emph{mutually consistent}.
This consistency manifests as a system of coboundary equations that the decoration data must satisfy.

\subsubsection{The 2+1D case}

Let us be explicit for $d=2$ spatial dimensions.
The decoration data is a triplet $(n_1, n_2, \nu_3)$ where:
\begin{itemize}
    \item $n_1 \in C^1(\Gb, \Zt)$ describes the Majorana chain decoration on 1-simplices (edges),
    \item $n_2 \in C^2(\Gb, \Zt)$ describes the complex fermion decoration on 2-simplices (triangles),
    \item $\nu_3 \in C^3(\Gb, \U(1))$ is the bosonic SPT $\U(1)$ phase factor on 3-simplices (tetrahedra in spacetime).
\end{itemize}

The consistency conditions form a hierarchical system of coboundary equations:
\begin{align}
\dd n_1 &= 0 \pmod{2}, \label{eq:2d-layer1} \\
\dd n_2 &= \omega_2 \cup n_1 + s_1 \cup n_1 \cup n_1 \pmod{2}, \label{eq:2d-layer2} \\
\dd_{s_1} \nu_3 &= \mathcal{O}_4[n_2]. \label{eq:2d-layer3}
\end{align}

Let us unpack these equations:

\paragraph{Equation \eqref{eq:2d-layer1}: Majorana chains must form closed loops.}
The Majorana chain decoration $n_1$ must be a cocycle (mod 2).
Physically, Majorana chains are 1D objects that must form closed loops on the domain walls---they cannot have dangling endpoints.
This is the simplest consistency condition.

\paragraph{Equation \eqref{eq:2d-layer2}: Complex fermion parity constraint.}
The coboundary of $n_2$ is not zero, but is determined by $n_1$ and the symmetry data $(\omega_2, s_1)$.
The right-hand side has two terms:
\begin{itemize}
    \item $\omega_2 \cup n_1$: the Majorana chains interact with the fermion parity extension,
    \item $s_1 \cup n_1 \cup n_1$: time-reversal symmetry produces additional Majorana modes that must be compensated by complex fermions.
\end{itemize}
This equation tells us that once we choose $n_1$, the complex fermion decoration $n_2$ is constrained: not every $n_2$ is allowed, only those whose failure to be a cocycle is exactly $\omega_2 \cup n_1 + s_1 \cup n_1 \cup n_1$.

\paragraph{Equation \eqref{eq:2d-layer3}: Bosonic phase obstruction.}
The twisted coboundary of $\nu_3$ is determined by a complicated obstruction function $\mathcal{O}_4[n_2]$, which involves cup products, higher cup products ($\cup_1$, $\cup_2$), and the symmetry data.
The explicit form is
\begin{align}
\mathcal{O}_4[n_2](01234) &= (-1)^{[\omega_2 \cup n_2 + n_2 \cup n_2 + n_2 \cup_1 \dd n_2 + \dd(s_1 \cup n_2 + n_2 \cup_2 \dd n_2)](01234)} \nonumber \\
&\quad \times (-1)^{\omega_2(013)\, \dd n_2(1234) + \dd n_2(0124)\, \dd n_2(0234)} \nonumber \\
&\quad \times (-i)^{\dd n_2(0123)\, [1 - \dd n_2(0124)] \pmod{2}}. \label{eq:O4-explicit}
\end{align}
We present this formula here not to derive it (which requires careful tracking of Majorana fermion phases under retriangulation), but to illustrate two important features:
\begin{enumerate}[label=(\roman*)]
    \item The obstruction function $\mathcal{O}_4$ depends on the \emph{lower layer} data $n_2$ (and through \eqref{eq:2d-layer2}, also on $n_1$).
    The layers are coupled: you cannot classify them independently.
    \item The formula is explicit and computable, involving only algebraic operations on cochains.
    This is what makes the entire classification amenable to computer algebra.
\end{enumerate}

\subsubsection{The 3+1D case}
\label{subsubsec:3d-case}

For $d=3$ spatial dimensions, the decoration data is $(n_1, n_2, n_3, \nu_4)$ with an additional layer:
\begin{itemize}
    \item $n_1 \in C^1(\Gb, \Z)$: the $p+ip$ superconductor decoration,
    \item $n_2 \in C^2(\Gb, \Zt)$: Majorana chain decoration,
    \item $n_3 \in C^3(\Gb, \Zt)$: complex fermion decoration,
    \item $\nu_4 \in C^4(\Gb, \U(1))$: bosonic SPT phase.
\end{itemize}
The consistency conditions follow the same hierarchical pattern:
\begin{align}
\dd n_1 &= 0, \label{eq:3d-layer1} \\
\dd n_2 &= n_2 \text{-dependent terms involving } \omega_2, s_1, n_1 \pmod{2}, \label{eq:3d-layer2} \\
\dd n_3 &= n_2 \cup n_2 + \omega_2 \cup n_2 + s_1 \cup (n_2 \cup_1 n_2) \pmod{2}, \label{eq:3d-layer3} \\
\dd \nu_4 &= \mathcal{O}_5[n_3, n_2], \label{eq:3d-layer4}
\end{align}
where the 3+1D obstruction function $\mathcal{O}_5$ has three components:
\begin{equation}
\mathcal{O}_5[n_3, n_2] = \mathcal{O}_5^c[n_3] \cdot \mathcal{O}_5^{c\gamma}[\dd n_3] \cdot \mathcal{O}_5^\gamma[n_2].
\end{equation}
These come from, respectively, the anticommutation of complex fermions, the cross terms between complex fermions and Majorana fermions, and the pure Majorana phase.

\subsection{The hierarchical structure: a preview of the spectral sequence}
\label{subsec:hierarchical-preview}

Let us step back and observe the structure of the coboundary equations.
There is a clear \textbf{hierarchy}:
\begin{enumerate}
    \item First solve for $n_1$: the equation $\dd n_1 = 0$ is independent of everything else.
    \item Then solve for $n_2$: the equation $\dd n_2 = F(n_1)$ depends on $n_1$ but not on $n_2$ or $\nu_3$.
    \item Finally solve for $\nu_3$: the equation $\dd \nu_3 = G(n_2, n_1)$ depends on both lower layers.
\end{enumerate}

This hierarchical structure is not accidental.
It reflects the fact that the decoration layers live on simplices of \emph{different codimensions}, and lower-codimensional data constrains higher-codimensional data.

Furthermore, at each level, the solving process has two aspects:
\begin{itemize}
    \item \textbf{Existence (obstruction)}: Given $n_1$, does a consistent $n_2$ exist? The right-hand side of \eqref{eq:2d-layer2} must itself be a cocycle (its coboundary must vanish). If it is not, then this particular $n_1$ is \emph{obstructed}---no consistent $n_2$ exists, and we must discard this $n_1$.
    \item \textbf{Uniqueness (equivalence)}: When solutions exist, different choices of $n_2$ may be equivalent if they differ by a coboundary of some lower-degree cochain. This reduces the space of distinct solutions.
\end{itemize}

This pattern of alternating obstruction and equivalence, applied layer by layer, is \emph{precisely} what a spectral sequence does.
In the next section, we will see how the Atiyah--Hirzebruch spectral sequence formalizes this intuitive picture into a rigorous and computable algebraic framework.


%% ============================================================
%% CHAPTER 3
%% ============================================================
\section{The Atiyah--Hirzebruch Spectral Sequence}
\label{sec:AHSS}

This is the mathematical core of the lecture.
We introduce the AHSS step by step, always keeping the physical decoration problem in sight.
The goal is not to present spectral sequences as abstract algebra, but to explain \emph{why they naturally arise} from the physics and \emph{how they help us compute}.

\subsection{Motivation: why we need more than group cohomology}
\label{subsec:motivation}

If the consistency equations in \S\ref{subsec:coboundary-eqs} were simply $\dd n_p = 0$ for every layer independently, the classification would reduce to a product of group cohomology groups:
\begin{equation}
\label{eq:naive-product}
\prod_{p} H^p(\Gb, A_{q(p)}).
\end{equation}
But the equations are \emph{coupled}: $\dd n_2$ depends on $n_1$, and $\dd \nu_3$ depends on $n_2$.
These couplings mean that:
\begin{itemize}
    \item Not every cocycle $n_1 \in Z^1(\Gb, \Zt)$ leads to a consistent decoration---some are obstructed by \eqref{eq:2d-layer2}.
    \item Two solutions $(n_1, n_2, \nu_3)$ and $(n_1, n_2', \nu_3')$ may be physically equivalent if they differ by a gauge transformation that mixes the layers.
\end{itemize}
The true classification is therefore \emph{smaller} than the naive product \eqref{eq:naive-product}, and the reduction is organized by the AHSS.

\subsection{Cochain complexes and cohomology: a brief reminder}
\label{subsec:cochain-reminder}

Before introducing spectral sequences, let us recall the basic objects.

\begin{definition}[Cochain complex]
A \emph{cochain complex} $(C^\bullet, \dd)$ is a sequence of abelian groups
\begin{equation}
\cdots \xrightarrow{\dd} C^{n-1} \xrightarrow{\dd} C^n \xrightarrow{\dd} C^{n+1} \xrightarrow{\dd} \cdots
\end{equation}
with $\dd \circ \dd = 0$.
Elements of $\ker(\dd: C^n \to C^{n+1})$ are called \emph{$n$-cocycles}, and elements of $\im(\dd: C^{n-1} \to C^n)$ are called \emph{$n$-coboundaries}.
\end{definition}

\begin{definition}[Cohomology]
The $n$-th \emph{cohomology group} is the quotient
\begin{equation}
H^n(C) := \frac{\ker(\dd: C^n \to C^{n+1})}{\im(\dd: C^{n-1} \to C^n)} = \frac{Z^n(C)}{B^n(C)}.
\end{equation}
\end{definition}

In our context, the cochain complex is built from the group $\Gb$ and a coefficient module $A$.
The group $C^p(\Gb, A)$ consists of all functions from $\Gb^p$ to $A$, and the coboundary $\dd$ is the standard group cohomology coboundary.

\subsection{Filtered complexes: the mathematical structure behind the layers}
\label{subsec:filtration}

The hierarchical structure of the decoration problem---different layers, coupled by cross-term obstructions---has a precise mathematical name: a \emph{filtered cochain complex}.

\begin{definition}[Decreasing filtration]
\label{def:filtration}
A \emph{decreasing filtration} on a cochain complex $(C^\bullet, \dd)$ is a sequence of subcomplexes
\begin{equation}
\label{eq:filtration}
\cdots \supseteq F^{p-1} C^n \supseteq F^p C^n \supseteq F^{p+1} C^n \supseteq \cdots \supseteq 0
\end{equation}
satisfying $\dd(F^p C^n) \subseteq F^p C^{n+1}$ (the coboundary respects the filtration).
\end{definition}

In the FSPT classification, the filtration has a direct physical interpretation.
Let the total cochain complex $C^n$ be the space of all decoration data of total degree $n = p + q$.
The filtration level $F^p C^n$ consists of decoration data where the first $p-1$ layers (the ones on \emph{low}-codimensional domain walls) are trivially zero:
\begin{equation}
\alpha \in F^p C^n \quad \Longleftrightarrow \quad \text{layers } 0, 1, \ldots, p-1 \text{ of } \alpha \text{ are zero.}
\end{equation}
The condition $\dd(F^p C^n) \subseteq F^p C^{n+1}$ holds because the coboundary of a cochain with vanishing low layers has vanishing low layers in its coboundary (the coboundary equations are ``upper-triangular'' in the layer index).

\begin{remark}
The filtration captures the idea that the layers have a natural ordering: you solve the bottom layer first, then use its solution to constrain the next layer, and so on.
A spectral sequence is the algebraic machine that implements this ``solve layer by layer'' strategy.
\end{remark}

\subsection{The bigrading: what $(p,q)$ means}
\label{subsec:bigrading}

The total degree $n$ of a cochain decomposes as $n = p + q$, where:
\begin{itemize}
    \item $p$ is the \emph{filtration degree}: it labels the layer (i.e., which codimensional domain wall the decoration lives on).
    Equivalently, $p$ is the degree of the inhomogeneous cochain on $\Gb$.
    \item $q$ is the \emph{complementary degree}: it labels the \emph{type} of the decorating topological phase.
\end{itemize}

For the FSPT classification, the $q$-index has a direct physical meaning:
\begin{equation}
\label{eq:q-dictionary}
\boxed{
\begin{aligned}
q = 0 &: \quad \text{Bosonic SPT phase} \quad (A_0 = \U(1)), \\
q = 1 &: \quad \text{Complex fermion} \quad (A_1 = \Zt), \\
q = 2 &: \quad \text{Majorana chain} \quad (A_2 = \Zt), \\
q = 3 &: \quad p+ip \text{ superconductor} \quad (A_3 = \Z).
\end{aligned}
}
\end{equation}

So for a $d$-dimensional spatial FSPT, the total degree is $n = d + 1$ (the spacetime dimension), and the relevant entries are those with $p + q = d + 1$, $p \geq 1$, $q \geq 0$.

\begin{example}[2+1D bigrading]
\label{ex:2d-bigrading}
For $d = 2$, the total degree is $n = 3$.
The relevant bidegrees are:
\begin{center}
\renewcommand{\arraystretch}{1.3}
\begin{tabular}{cccl}
\toprule
$(p,q)$ & $p$-cochain on $\Gb$ & Coefficient $A_q$ & Physical layer \\
\midrule
$(1,2)$ & 1-cochain & $\Zt$ & Majorana chain \\
$(2,1)$ & 2-cochain & $\Zt$ & Complex fermion \\
$(3,0)$ & 3-cochain & $\U(1)$ & Bosonic SPT \\
\bottomrule
\end{tabular}
\end{center}
This is exactly the triplet $(n_1, n_2, \nu_3)$ from \S\ref{subsec:coboundary-eqs}.
\end{example}

\begin{example}[3+1D bigrading]
For $d = 3$, the total degree is $n = 4$.
The relevant bidegrees are:
\begin{center}
\renewcommand{\arraystretch}{1.3}
\begin{tabular}{ccl}
\toprule
$(p,q)$ & Coefficient $A_q$ & Physical layer \\
\midrule
$(1,3)$ & $\Z$ & $p+ip$ superconductor \\
$(2,2)$ & $\Zt$ & Majorana chain \\
$(3,1)$ & $\Zt$ & Complex fermion \\
$(4,0)$ & $\U(1)$ & Bosonic SPT \\
\bottomrule
\end{tabular}
\end{center}
\end{example}

\subsection{The $E_2$ page: the first approximation}
\label{subsec:E2-page}

Given the bigrading, the simplest thing we can do is to classify each layer \emph{independently}, ignoring all cross-layer couplings.
This means computing the group cohomology of $\Gb$ with each coefficient module separately:
\begin{equation}
\label{eq:E2-page}
E_2^{p,q} := H^p(\Gb, A_q).
\end{equation}

\begin{definition}[$E_2$ page]
The \emph{second page} (or $E_2$ page) of the AHSS is the collection of all groups $E_2^{p,q}$, arranged in a grid with $p$ on the horizontal axis and $q$ on the vertical axis.
\end{definition}

Why is this called the ``second page'' and not the ``first''?
The first page $E_1^{p,q} = C^p(\Gb, A_q)$ is the full cochain module, before taking any cohomology.
The $E_2$ page takes the cohomology of the standard group coboundary, which is a well-defined operation independent of the cross-layer couplings.
It is the natural starting point for classification.

\begin{remark}[Physical interpretation of $E_2$]
\label{rem:E2-physical}
$E_2^{p,q} = H^p(\Gb, A_q)$ is the group of \emph{inequivalent cocycles} in degree $p$ with coefficient $A_q$, modulo coboundaries.
Physically, this is the classification of $A_q$-type decorations on $p$-dimensional domain walls, if those decorations were completely independent of all other layers.
The $E_2$ page is the ``zeroth-order'' classification: it tells you how many candidate decorations there are, before accounting for inter-layer obstructions.
\end{remark}

\begin{example}[2+1D $E_2$ page for $\Gb = \Zt$]
For $\Gb = \Zt$ (e.g., time-reversal symmetry alone), the group cohomology is $H^n(\Zt, \Zt) = \Zt$ for all $n \geq 0$, and similarly for $\U(1)$ coefficients.
So the $E_2$ page along the line $p + q = 3$ is:
\begin{equation}
E_2^{1,2} = H^1(\Zt, \Zt) = \Zt, \quad
E_2^{2,1} = H^2(\Zt, \Zt) = \Zt, \quad
E_2^{3,0} = H^3(\Zt, \U(1)) = \Zt.
\end{equation}
If there were no cross-layer couplings, the classification would be $\Zt^3$.
As we will see, the differentials may reduce this.
\end{example}

\subsection{Differentials: obstructions from cross-layer couplings}
\label{subsec:differentials}

The $E_2$ page ignores the couplings between layers.
The \emph{differentials} of the spectral sequence encode precisely these couplings.

\begin{definition}[Differential]
The $r$-th differential is a homomorphism
\begin{equation}
\label{eq:dr-bidegree}
d_r: E_r^{p,q} \longrightarrow E_r^{p+r,\, q-r+1}
\end{equation}
satisfying $d_r \circ d_r = 0$.
\end{definition}

The bidegree shift $(+r, -r+1)$ is universal for all cohomological spectral sequences.
Let us see why this particular shift arises from physics.

\paragraph{Physical interpretation of $d_2$.}
The first nontrivial differential is $d_2: E_2^{p,q} \to E_2^{p+2, q-1}$.
In the 2+1D case, this gives
\begin{equation}
d_2: E_2^{1,2} \longrightarrow E_2^{3,1}.
\end{equation}
This maps the Majorana layer $(p,q) = (1,2)$ to the complex fermion layer $(p,q) = (3,1)$.
Wait---$(3,1)$ has $p + q = 4 \neq 3$.
The total degree has changed by $+1$!
This is correct: the differential $d_2$ raises the total degree by 1, reflecting the fact that the \emph{obstruction} to extending a solution from one layer to the next lives one degree higher.

Let us connect this to the coboundary equations.
Consider a Majorana cocycle $n_1 \in Z^1(\Gb, \Zt)$, representing an element of $E_2^{1,2}$.
The coboundary equation \eqref{eq:2d-layer2} says
\begin{equation}
\dd n_2 = \omega_2 \cup n_1 + s_1 \cup n_1 \cup n_1 \pmod{2}.
\end{equation}
For this equation to have a solution $n_2$, the right-hand side must itself be a coboundary (i.e., it must be cohomologically trivial).
The map
\begin{equation}
d_2(n_1) := [\omega_2 \cup n_1 + s_1 \cup n_1 \cup n_1] \in H^3(\Gb, \Zt) = E_2^{3,1}
\end{equation}
is precisely the obstruction.
If $d_2(n_1) \neq 0$, then this particular $n_1$ is \emph{obstructed}: there is no consistent complex fermion decoration compatible with it.

\begin{remark}
More precisely, $d_2$ is not just defined on cocycles $n_1$ but on cohomology classes $[n_1] \in H^1(\Gb, \Zt)$.
One can verify that the obstruction depends only on the cohomology class, not on the particular cocycle representative.
This is a standard check in the construction of spectral sequence differentials.
\end{remark}

\paragraph{Physical interpretation of higher differentials.}
The $d_3$ differential connects $E_3^{p,q}$ to $E_3^{p+3, q-2}$, skipping two bidegree positions.
Physically, this represents a \emph{secondary obstruction}: even after resolving the $d_2$ obstruction, there may be a further obstruction that only becomes visible after the first-order correction is made.

In the FSPT context, $d_3$ and higher differentials typically encode the obstruction of the bosonic phase $\nu_{d+1}$ by the fermionic decoration data.
For example, the complicated obstruction function $\mathcal{O}_4[n_2]$ in \eqref{eq:O4-explicit} contributes to both $d_2$ (through its dependence on $n_2$, which in turn depends on $n_1$) and potentially $d_3$ (through higher-order effects).

\subsection{Pages: the ``survivors'' after each round}
\label{subsec:pages}

After defining $d_r$, the next page is obtained by taking cohomology:

\begin{definition}[Page recursion]
\begin{equation}
\label{eq:page-recursion}
E_{r+1}^{p,q} := \frac{\ker\bigl(d_r: E_r^{p,q} \to E_r^{p+r,\, q-r+1}\bigr)}{\im\bigl(d_r: E_r^{p-r,\, q+r-1} \to E_r^{p,q}\bigr)}.
\end{equation}
\end{definition}

Let us parse this formula carefully.

The \textbf{numerator} $\ker(d_r)$ is the subgroup of $E_r^{p,q}$ consisting of elements that are \emph{not obstructed} by $d_r$.
These are the ``survivors'' that pass the obstruction test.

The \textbf{denominator} $\im(d_r)$ is the subgroup of $E_r^{p,q}$ consisting of elements that are ``hit'' by $d_r$ from another bidegree.
Being hit by a differential means being the image of an obstruction, which corresponds to a \emph{redundancy}: these elements can be trivialized by a gauge transformation.

\begin{remark}[An elimination tournament]
The progression $E_2 \to E_3 \to E_4 \to \cdots$ can be thought of as an elimination tournament:
\begin{itemize}
    \item $E_2$: all candidates (group cohomology classes).
    \item $d_2$ eliminates obstructed candidates and identifies redundant ones.
    \item $E_3$: survivors of the first round.
    \item $d_3$ eliminates further obstructed candidates.
    \item $E_4$: survivors of the second round.
    \item And so on.
\end{itemize}
At each round, the groups can only get \emph{smaller} (or stay the same).
The process terminates when no further differentials can act.
\end{remark}

\subsection{Convergence: the $E_\infty$ page}
\label{subsec:convergence}

For bounded spectral sequences (which is always the case in our setting, since $p \geq 0$ and $q \geq 0$), the pages eventually stabilize:
\begin{equation}
\label{eq:stabilization}
E_r^{p,q} = E_{r+1}^{p,q} = \cdots = E_\infty^{p,q} \qquad \text{for } r \gg 0.
\end{equation}

\begin{definition}[Convergence]
The spectral sequence \emph{converges} to the cohomology of the total complex:
\begin{equation}
\label{eq:convergence}
E_\infty^{p,q} \cong \gr^p H^{p+q}(C),
\end{equation}
where $\gr^p H^n(C) := F^p H^n(C) / F^{p+1} H^n(C)$ is the associated graded of the filtration on cohomology.
\end{definition}

What does this mean physically?
The total cohomology $H^{d+1}(C)$ is the \emph{true classification group} of FSPT phases in $d$ spatial dimensions.
The $E_\infty$ page gives the \emph{graded pieces} of this group.
Each graded piece $E_\infty^{p,q}$ tells us how many independent phases survive in the layer with bidegree $(p,q)$, after all obstructions and equivalences have been accounted for.

\begin{remark}[Warning: $E_\infty$ is not the whole answer]
\label{rem:Einfty-warning}
The $E_\infty$ page gives the \emph{graded pieces}, not the classification group itself.
Knowing the graded pieces $\gr^p H^n$ does not uniquely determine $H^n$ as an abelian group.
The passage from graded pieces to the full group requires solving the \emph{extension problem}, which we will discuss in detail in Part II (\S\ref{sec:extension}).

As a simple analogy: knowing that a group has composition factors $\Zt$ and $\Zt$ does not tell you whether the group is $\Zt \times \Zt$ or $\Zn{4}$.
\end{remark}


\subsection{The AHSS in action: a complete 2+1D example}
\label{subsec:AHSS-example}

Let us work through the full AHSS for 2+1D FSPT with $\Gb = \Zt \times \Zt$ (two independent $\Zt$ symmetries, both unitary), with trivial $\omega_2 = 0$ and $s_1 = 0$.

\paragraph{Step 1: Compute $E_2$.}
The group cohomology of $\Gb = \Zt \times \Zt$ with $\Zt$ coefficients is computed by the Künneth formula.
The relevant entries at total degree $n = 3$ are:
\begin{align}
E_2^{1,2} &= H^1(\Zt \times \Zt, \Zt) = \Zt^2, \\
E_2^{2,1} &= H^2(\Zt \times \Zt, \Zt) = \Zt^3, \\
E_2^{3,0} &= H^3(\Zt \times \Zt, \U(1)) = \Zt^4.
\end{align}

\paragraph{Step 2: Compute $d_2$.}
Since $\omega_2 = 0$ and $s_1 = 0$, the differential $d_2: E_2^{1,2} \to E_2^{3,1}$ is the map
$d_2(n_1) = [\omega_2 \cup n_1 + s_1 \cup n_1 \cup n_1] = 0.$
So $d_2 = 0$ in this case, and no Majorana cocycles are obstructed.

\paragraph{Step 3: Compute $E_3$.}
Since $d_2 = 0$ everywhere along the $n = 3$ line, we have $E_3^{p,q} = E_2^{p,q}$ for all $(p,q)$ with $p + q = 3$.

\paragraph{Step 4: Check $d_3$ and higher.}
For this particular symmetry, the higher differentials also turn out to be zero, so $E_\infty = E_2$ along the line $p + q = 3$.
The classification (before solving the extension problem) has graded pieces:
\begin{equation}
E_\infty^{1,2} = \Zt^2, \quad E_\infty^{2,1} = \Zt^3, \quad E_\infty^{3,0} = \Zt^4.
\end{equation}

The naive direct sum would give $\Zt^9$.
But whether the true classification group is $\Zt^9$ or something else (e.g., $\Zn{4}$ appearing) depends on the extension problem, which we will address in Part II.

\subsection{Dictionary: physics $\leftrightarrow$ AHSS}
\label{subsec:dictionary}

We now summarize the complete correspondence between the physical domain wall decoration and the AHSS:

\begin{center}
\renewcommand{\arraystretch}{1.5}
\begin{tabular}{>{\raggedright}p{6.5cm} p{6.5cm}}
\toprule
\textbf{Physics} & \textbf{AHSS} \\
\midrule
Decoration layers on domain walls of codimension $q$
& Bigrading $(p,q)$: the coefficient $A_q$ labels the type of decoration, $p$ labels the cochain degree \\
\midrule
Candidate decorations (ignoring inter-layer coupling)
& $E_2^{p,q} = H^p(\Gb, A_q)$ \\
\midrule
Consistency equations coupling different layers (e.g., $\dd n_2 = F(n_1)$)
& Differentials $d_r: E_r^{p,q} \to E_r^{p+r, q-r+1}$ \\
\midrule
A cocycle $n_1$ is obstructed: no consistent $n_2$ exists
& $d_2([n_1]) \neq 0$ in $E_2^{p+2, q-1}$ \\
\midrule
Two solutions $(n_1, n_2, \nu_3)$ differ by a gauge transformation
& Quotient by $\im(d_r)$ in the page recursion \\
\midrule
``Survivors'' after accounting for all obstructions and equivalences
& $E_\infty^{p,q}$: the stabilized page \\
\midrule
Total classification of FSPT phases in $d$ dimensions
& $H^{d+1}(C)$, reconstructed from $E_\infty^{p,q}$ via extensions \\
\bottomrule
\end{tabular}
\end{center}

\subsection{Why AHSS is more than bookkeeping}
\label{subsec:AHSS-power}

One might ask: if the coboundary equations already encode all the physics, why bother with the spectral sequence?
The answer is that AHSS provides three things that working directly with coboundary equations does not:

\begin{enumerate}[label=(\roman*)]
    \item \textbf{Systematic organization.}
    For complicated symmetry groups (e.g., space groups with hundreds of generators), the coboundary equations form a massive coupled system.
    The spectral sequence breaks this into a sequence of \emph{independent} linear algebra problems on each page, which can be solved one at a time.

    \item \textbf{Computability.}
    Each page $E_r$ is obtained from $E_{r-1}$ by a standard algebraic operation (taking cohomology of a linear map).
    This is amenable to computer algebra.
    The \texttt{SptSet} package exploits this by implementing the page recursion directly.

    \item \textbf{Structural insight.}
    The spectral sequence reveals which obstructions are ``primary'' (visible at $E_2$), which are ``secondary'' (only appearing at $E_3$), and so on.
    This hierarchy has physical meaning: primary obstructions come from simple consistency conditions, while secondary obstructions arise from subtler interference effects.
\end{enumerate}


%% ============================================================
%% CHAPTER 4
%% ============================================================
\section{Computing Classification with SptSet}
\label{sec:sptset-classification}

The mathematical framework of the AHSS has been implemented as the \texttt{SptSet} package for the GAP computer algebra system.
This section gives a hands-on introduction: how to set up a computation, how to read the output, and what the output means mathematically and physically.

\subsection{The computational pipeline}
\label{subsec:pipeline}

The standard workflow for computing an FSPT classification with \texttt{SptSet} consists of five steps:

\begin{equation}
\boxed{\text{Group}} \;\xrightarrow{1}\; \boxed{\text{Resolution}} \;\xrightarrow{2}\; \boxed{\text{Parity map}} \;\xrightarrow{3}\; \boxed{\text{Spectral sequence}} \;\xrightarrow{4}\; \boxed{\text{Classification}}
\end{equation}

Let us walk through each step using a concrete 2D example.

\subsubsection{Step 1: Choose a symmetry group}

For crystalline symmetries, the symmetry group $\Gb$ is a space group.
In GAP, 2D space groups (wallpaper groups) are created using the \texttt{cryst} package:
\begin{verbatim}
LoadPackage("SptSet");
SG := SpaceGroupBBNWZ(2, 13);   # wallpaper group No. 13 (p6)
\end{verbatim}
Here \texttt{SpaceGroupBBNWZ(2, it)} creates the \texttt{it}-th wallpaper group ($\texttt{it} = 2, \ldots, 17$).
The space group must be converted to a polycyclically presented (Pcp) group, which is the format required by the HAP package for computing resolutions:
\begin{verbatim}
fSG := IsomorphismPcpGroup(SG);
SG1 := Image(fSG);
\end{verbatim}

\subsubsection{Step 2: Compute a resolution}

A resolution of the group is the algebraic input needed for group cohomology computations.
For crystallographic groups, HAP provides:
\begin{verbatim}
R := ResolutionAlmostCrystalGroup(SG1, 6);
\end{verbatim}
The second argument (\texttt{6}) is the dimension bound: the resolution is computed up to degree 6, which is sufficient for classifying FSPT up to 4+1D ($p + q \leq 5$).

\subsubsection{Step 3: Define the fermion parity map}

The fermion parity map $\rho: \Gb \to \text{GL}(1, \Z) \cong \{\pm 1\}$ encodes whether each symmetry element is unitary or anti-unitary.
For space groups, this is determined by the determinant of the point-group part:
\begin{verbatim}
gs := GeneratorsOfGroup(SG);
f := GroupHomomorphismByImagesNC(SG1, GL(1, Integers),
  List(gs, x -> Image(fSG, x)),
  List(gs, x -> [[DeterminantMat(x)]]));
\end{verbatim}
The function $s(g) = \frac{1 - \rho(g)_{11}}{2}$ then gives $s(g) = 0$ for unitary and $s(g) = 1$ for anti-unitary elements.

\subsubsection{Step 4: Construct the spectral sequence}

With the resolution and the parity map in hand, we construct the spectral sequence object:
\begin{verbatim}
SS := FermionEZSPTSpecSeq(R, f);
\end{verbatim}
This single function call does a great deal of work behind the scenes:
\begin{enumerate}[label=(\alph*)]
    \item It sets up the spectrum of coefficient systems $\{A_q\}$: $\U(1)$, $\Zt$, $\Zt$, $\Z$.
    \item It installs all the cochain-level coboundary formulas (the physics input) for each bidegree $(r, p, q)$.
    \item It installs the stacking twisters (for later use in the extension problem).
    \item It initializes the caching infrastructure for pages $E_r^{p,q}$.
\end{enumerate}

\subsubsection{Step 5: Read off the classification}

The classification layers are displayed by:
\begin{verbatim}
FermionSPTLayersVerbose(SS, 2);
\end{verbatim}
The argument \texttt{2} is the spatial dimension $d$.
The function prints the spectral sequence pages $E_r^{p,q}$ for all relevant bidegrees along $p + q = d + 1 = 3$.


\subsection{Reading the output}
\label{subsec:reading-output}

A typical output looks like:
\begin{verbatim}
Majorana:       # (p,q) = (1,2)
  E2: <ZL-Module with torsions [ 2 ]>
  E3: <ZL-Module with torsions [ 2 ]>
  E4: <ZL-Module with torsions [ 2 ]>
Complex fermion:  # (p,q) = (2,1)
  E2: <ZL-Module with torsions [ 2, 2 ]>
  E3: <ZL-Module with torsions [ 2 ]>
Bosonic:        # (p,q) = (3,0)
  E2: <ZL-Module with torsions [ 2, 2, 2 ]>
  E3: <ZL-Module with torsions [ 2, 2, 2 ]>
\end{verbatim}

Here is how to read it:
\begin{itemize}
    \item Each heading (``Majorana'', ``Complex fermion'', ``Bosonic'') corresponds to a value of $q$ according to the dictionary \eqref{eq:q-dictionary}.
    \item Under each heading, the lines show the groups $E_r^{p,q}$ for $r = 2, 3, \ldots$
    \item \texttt{<ZL-Module with torsions [2, 2]>} means the group $\Zt \times \Zt$.
    \item When the group stops changing as $r$ increases, the spectral sequence has \emph{stabilized} at that bidegree: the displayed group is $E_\infty^{p,q}$.
\end{itemize}

In the example above:
\begin{itemize}
    \item The Majorana layer $E_\infty^{1,2} = \Zt$.
    \item The complex fermion layer changed from $\Zt^2$ at $E_2$ to $\Zt$ at $E_3$ (meaning a nontrivial $d_2$ killed one generator).
    \item The bosonic layer $E_\infty^{3,0} = \Zt^3$.
\end{itemize}

\subsection{Complete example script}
\label{subsec:complete-script}

Here is the complete script that classifies all 2D wallpaper groups:
\begin{verbatim}
LoadPackage("SptSet");

for it in [2..17] do
  Display(it);
  SG := SpaceGroupBBNWZ(2, it);
  fSG := IsomorphismPcpGroup(SG);
  SG1 := Image(fSG);
  R := ResolutionAlmostCrystalGroup(SG1, 6);
  gs := GeneratorsOfGroup(SG);

  f := GroupHomomorphismByImagesNC(SG1, GL(1, Integers),
    List(gs, x -> Image(fSG, x)),
    List(gs, x -> [[DeterminantMat(x)]]));

  SS := FermionEZSPTSpecSeq(R, f);
  FermionSPTLayersVerbose(SS, 2);
od;
\end{verbatim}
This loops over all 16 non-trivial wallpaper groups, computing the AHSS for each one.
The output gives the $E_\infty$ pages (graded pieces of the classification) for all 2+1D FSPT phases protected by each wallpaper group symmetry.

\subsection{Architecture of SptSet: a map for the curious}
\label{subsec:architecture}

For readers who want to look under the hood, here is a map of how the code is organized:

\begin{center}
\renewcommand{\arraystretch}{1.3}
\begin{tabular}{ll}
\toprule
\textbf{File} & \textbf{Role} \\
\midrule
\texttt{lib/ss\_fermion\_ez.gi} & Physics input: coboundary formulas and twisters \\
\texttt{lib/ss\_vanilla.gi} & AHSS engine: pages, differentials, caching \\
\texttt{lib/spectral\_sequence.gi} & Spectral sequence algebra (page recursion) \\
\texttt{lib/bar\_resolution\_map*.gi} & Bridge between HAP resolutions and cochains \\
\texttt{lib/group\_cohomology.gi} & Group cohomology via HAP \\
\texttt{lib/coefficient.gi} & Coefficient systems ($\U(1)$, $\Zt$, etc.) \\
\texttt{lib/zlmap.gi} & Linear algebra over $\Z$ \\
\bottomrule
\end{tabular}
\end{center}

The separation between physics input (\texttt{ss\_fermion\_ez.gi}) and the general AHSS engine (\texttt{ss\_vanilla.gi}) is deliberate.
To extend the classification to new physical models, one writes a new physics-input file (analogous to \texttt{ss\_fermion\_ez.gi}) that installs the appropriate coboundary formulas, while the AHSS engine remains unchanged.

The two key APIs that connect physics to the engine are:
\begin{itemize}
    \item \texttt{SptSetInstallCoboundary(SS, r, p, q, f)}: installs a cochain-level formula $f$ that defines the $d_r$ differential at bidegree $(p,q)$.
    \item \texttt{SptSetInstallAddTwister(SS, p, q, A)}: installs a stacking correction (twister) at bidegree $(p,q)$, used for the extension problem.
\end{itemize}
Both take functions as arguments: the physics formulas are literally GAP functions that take cochains in and return cochains out.


%% ============================================================
%% CHAPTER 5
%% ============================================================
\section{Stacking: Physics of the Group Law}
\label{sec:stacking-physics}

Part I answered the question: \emph{what are the distinct FSPT phases?}
Part II addresses the follow-up: \emph{how do these phases combine when you stack two systems together?}

\subsection{What is stacking?}
\label{subsec:what-is-stacking}

Given two FSPT states $\Psi_1$ and $\Psi_2$ with the same symmetry $\Gf$, their \emph{stacking} (or \emph{tensor product}) $\Psi_1 \boxtimes \Psi_2$ is the state obtained by placing the two systems on top of each other and treating them as a single system.

The set of FSPT phases forms an \emph{abelian group} under stacking:
\begin{itemize}
    \item The identity is the trivial (product) state.
    \item The inverse of a phase is its ``mirror image'' (orientation reversal or complex conjugation).
    \item Stacking is associative and commutative.
\end{itemize}
The goal of this section is to determine this group---the \emph{stacking group}---which is the true classification group of FSPT phases.

\subsection{Naive expectation vs.\ reality}
\label{subsec:naive-vs-reality}

If each layer of the decoration were independent, the stacking rule would be simple: add each layer separately.
Given two decorations $(n_1, n_2, \nu_3)$ and $(n_1', n_2', \nu_3')$, their stack would be
\begin{equation}
(n_1 + n_1',\; n_2 + n_2',\; \nu_3 \cdot \nu_3').
\end{equation}
This would make the stacking group a direct product: $E_\infty^{1,2} \times E_\infty^{2,1} \times E_\infty^{3,0}$.

However, the layers are \emph{not} independent.
The physical stacking rules contain \emph{cross-terms} between different layers.
These cross-terms arise because the decoration on a domain wall depends on the configuration of the \emph{entire} system, including the decorations on other domain walls.

\subsection{The 2+1D stacking rules}
\label{subsec:2d-stacking}

The stacking of two 2+1D FSPT states $(n_1, n_2, \nu_3) \boxtimes (n_1', n_2', \nu_3') = (N_1, N_2, \mathcal{V}_3)$ is given by:
\begin{align}
N_1 &= n_1 + n_1' \pmod{2}, \label{eq:stack-N1} \\
N_2 &= n_2 + n_2' + \underbrace{n_1 \cup n_1' + s_1 \cup (n_1 \cup_1 n_1')}_{m_2} \pmod{2}, \label{eq:stack-N2} \\
\mathcal{V}_3 &= \nu_3 \cdot \nu_3' \cdot \underbrace{(-1)^{\epsilon_3} \cdot e^{2\pi i \theta_3[n_1, n_1']}}_{\mathcal{E}_3}. \label{eq:stack-V3}
\end{align}

Let us examine each equation:

\paragraph{Equation \eqref{eq:stack-N1}: The Majorana layer adds simply.}
The Majorana chains combine by simple mod-2 addition.
No cross-terms appear at this level.
This is because the Majorana layer is the ``lowest'' layer (highest $q$), and there is nothing below it to couple to.

\paragraph{Equation \eqref{eq:stack-N2}: Cross-terms in the complex fermion layer.}
The complex fermion decoration $N_2$ is \emph{not} simply $n_2 + n_2'$.
There is an extra term $m_2 = n_1 \cup n_1' + s_1 \cup (n_1 \cup_1 n_1')$ that depends on the \emph{Majorana layers} of both states.
Physically, when two Majorana chain configurations overlap, additional complex fermion modes can be created at their intersections.
These must be properly accounted for in $N_2$.

\paragraph{Equation \eqref{eq:stack-V3}: Cross-terms in the bosonic layer.}
The bosonic phase $\mathcal{V}_3$ picks up a correction factor $\mathcal{E}_3$ that depends on all lower layers.
This factor has two parts:
\begin{itemize}
    \item $(-1)^{\epsilon_3}$: a sign from the anticommutation of complex fermions, with $\epsilon_3$ involving cup products and higher cup products of $n_2, n_2', N_2, m_2$ and their coboundaries.
    \item $e^{2\pi i \theta_3}$: a phase from the recombination of Majorana fermions during stacking, the so-called ``Majorana phase''.
    This is a $\U(1)$-valued function of $n_1$ and $n_1'$ that involves $1/2$, $1/4$, and even $5/8$ fractional coefficients in the exponent, reflecting the $\Zn{8}$-valued nature of certain Majorana invariants.
\end{itemize}

\subsection{The key observation: stacking is not a direct product}
\label{subsec:key-observation}

The cross-terms $m_2$ in \eqref{eq:stack-N2} and $\mathcal{E}_3$ in \eqref{eq:stack-V3} have a profound consequence: the stacking group is \emph{not} a direct product of the individual layer classifications.

\begin{example}
Consider the $\Gf = \Zn{4}^{f,T}$ symmetry in 3+1D.
The spectral sequence gives:
\begin{equation}
E_\infty^{1,3} = \Zt \;\; (p+ip), \qquad E_\infty^{3,1} = \Zt \;\; (\text{complex fermion}).
\end{equation}
If stacking were a direct product, the classification group would be $\Zt \times \Zt = \Zn{2}^2$.
But the cross-terms in the stacking rule produce a nontrivial extension:
the $p+ip$ generator, when stacked with itself, produces a nonzero complex fermion class.
In other words, $2 \times [p+ip] = [\text{CF}]$ in the full classification.
This means the true classification group is $\Zn{4}$, not $\Zn{2}^2$.
\end{example}

This is the essence of the \emph{extension problem}, which we formalize mathematically in the next section.


%% ============================================================
%% CHAPTER 6
%% ============================================================
\section{The Extension Problem: Mathematics of Stacking}
\label{sec:extension}

The AHSS gives us the graded pieces $E_\infty^{p,q}$.
But the true classification group $\mathcal{C}$ is generally \emph{not} the direct sum of these pieces.
Reconstructing $\mathcal{C}$ from its graded pieces is the \emph{extension problem}.
This section explains the mathematics behind this problem and how it connects to the physical stacking rules.

\subsection{The extension problem: a general formulation}
\label{subsec:extension-general}

\begin{definition}[Extension of abelian groups]
An \emph{extension} of an abelian group $B$ by an abelian group $A$ is a short exact sequence
\begin{equation}
\label{eq:ses}
0 \longrightarrow A \xrightarrow{\iota} E \xrightarrow{\pi} B \longrightarrow 0.
\end{equation}
The group $E$ ``extends'' $B$ by $A$: it contains $A$ as a subgroup, with quotient $B$.
\end{definition}

The extension is \emph{trivial} (or \emph{splits}) if $E \cong A \oplus B$.
In general, there can be nontrivial extensions.

\begin{example}
The groups $\Zt \times \Zt$ and $\Zn{4}$ are both extensions of $\Zt$ by $\Zt$:
\begin{align}
0 \to \Zt \to \Zt \times \Zt \to \Zt \to 0 &\qquad \text{(trivial extension)}, \\
0 \to \Zt \to \Zn{4} \to \Zt \to 0 &\qquad \text{(nontrivial extension)}.
\end{align}
In the first case, $1 + 1 = 0$ in each factor independently.
In the second case, the generator of $\Zn{4}$ has order 4: $1 + 1 = 2 \neq 0$, but $2 + 2 = 0$.
The element $2 \in \Zn{4}$ is the nontrivial element of the subgroup $\Zt$.
\end{example}

\subsection{Extension cocycles: the mathematical incarnation of cross-terms}
\label{subsec:extension-cocycles}

Given a short exact sequence \eqref{eq:ses}, choose a set-theoretic section $s: B \to E$ (a map with $\pi \circ s = \mathrm{id}_B$, not necessarily a group homomorphism).
The failure of $s$ to be a homomorphism is measured by the \emph{extension 2-cocycle}:
\begin{equation}
\label{eq:ext-cocycle}
t(b_1, b_2) := s(b_1) + s(b_2) - s(b_1 + b_2) \in A.
\end{equation}
The map $t: B \times B \to A$ satisfies the cocycle condition $\delta t = 0$ (associativity of $E$ forces this).
The extension class $[t] \in H^2(B, A)$ determines the isomorphism type of $E$ up to equivalence.

\begin{remark}[Connection to physics]
Equation \eqref{eq:ext-cocycle} has exactly the same structure as the stacking cross-terms.
If we identify:
\begin{align*}
B &= E_\infty^{p_1, q_1} \quad (\text{one layer}), \\
A &= E_\infty^{p_2, q_2} \quad (\text{another layer, with } p_2 > p_1), \\
s(b) &= \text{the cochain representative of the class } b, \\
t(b_1, b_2) &= \text{the stacking cross-term (add-twister)},
\end{align*}
then the extension cocycle \eqref{eq:ext-cocycle} is precisely the \texttt{addTwister} function installed in SptSet.
\end{remark}

\subsection{From graded pieces to the full classification}
\label{subsec:graded-to-full}

For a total degree $n$, the AHSS produces a filtration
\begin{equation}
0 = F_{-1} \subset F_0 \subset F_1 \subset \cdots \subset F_n = \mathcal{C}_n,
\end{equation}
with graded pieces
\begin{equation}
F_p / F_{p-1} \cong E_\infty^{p, n-p}.
\end{equation}
The full classification group $\mathcal{C}_n$ is reconstructed by solving a chain of extension problems:
\begin{equation}
\label{eq:extension-chain}
\begin{tikzcd}[column sep=small]
0 \ar[r] & E_\infty^{1, n-1} \ar[r] & F_1 \ar[r] & E_\infty^{0, n} \ar[r] & 0 \\
0 \ar[r] & F_1 \ar[r] & F_2 \ar[r] & E_\infty^{2, n-2} \ar[r] & 0 \\
& & \vdots & & \\
0 \ar[r] & F_{n-1} \ar[r] & \mathcal{C}_n \ar[r] & E_\infty^{n, 0} \ar[r] & 0
\end{tikzcd}
\end{equation}

Each row is an extension problem that asks: given $F_{p-1}$ and $E_\infty^{p, n-p}$, what is $F_p$?
The answer is determined by the extension cocycle, which in our setting is computed from the stacking cross-terms.

\begin{remark}
In many cases (especially for odd-order groups or when the cross-terms happen to be trivial), the extensions split and $\mathcal{C}_n \cong \bigoplus_p E_\infty^{p, n-p}$.
The interesting physics happens when extensions are nontrivial, leading to groups like $\Zn{4}$ or $\Zn{8}$ instead of products of $\Zt$'s.
\end{remark}

\subsection{Add-twisters in SptSet: the computational implementation}
\label{subsec:addtwister-code}

In the \texttt{SptSet} package, the stacking of two cochain representatives is implemented by the function \texttt{SptSetStack(c1, c2)}.
For each layer $p$ (with $q = n - p$), the stacking rule is:
\begin{equation}
\label{eq:stack-code}
\texttt{layers[p+1]} := \texttt{c1.layers[p+1]} + \texttt{c2.layers[p+1]} + \texttt{addTwister[p+1][q+1](c1.layers, c2.layers)}.
\end{equation}
This is equation \eqref{eq:ext-cocycle} in code form: the first two terms are the ``naive'' addition, and the third term is the cross-term $t(b_1, b_2)$.

The twister function \texttt{addTwister[p+1][q+1]} takes as input the \emph{full layer lists} of both cochains, not just the individual layers.
This is because the cross-term for a given $(p,q)$ can depend on layers at \emph{different} bidegrees.
For example, the bosonic twister at $(p,q) = (3,0)$ in 2+1D depends on the Majorana layers $(p,q) = (1,2)$ and complex fermion layers $(p,q) = (2,1)$ of both inputs.

\subsection{The worked $\Zn{4}^{f,T}$ example revisited}
\label{subsec:z4ft-revisited}

Let us trace through the extension problem for the $\Gf = \Zn{4}^{f,T}$ example in 3+1D.

\paragraph{Step 1: The graded pieces.}
The AHSS gives:
\begin{equation}
E_\infty^{1,3} = \Zt \quad (p+ip), \qquad E_\infty^{2,2} = 0, \qquad E_\infty^{3,1} = \Zt \quad (\text{CF}), \qquad E_\infty^{4,0} = 0.
\end{equation}
The Majorana and bosonic layers vanish after the differentials act.

\paragraph{Step 2: Extension chain.}
We need to solve:
\begin{equation}
0 \to E_\infty^{3,1} \to \mathcal{C}_4 \to E_\infty^{1,3} \to 0,
\end{equation}
i.e., $0 \to \Zt \to \mathcal{C}_4 \to \Zt \to 0$.
The two possibilities are $\mathcal{C}_4 = \Zt \times \Zt$ (trivial extension) or $\mathcal{C}_4 = \Zn{4}$ (nontrivial).

\paragraph{Step 3: Compute the extension cocycle.}
Let $x$ be the generator of $E_\infty^{1,3}$ (the $p+ip$ class).
We stack $x$ with itself: $x \boxtimes x = 2x$.
At the cochain level, the stacking produces a nonzero term in the complex fermion layer, through the add-twister.
After purification (removing coboundary-trivial parts), the $p = 1$ layer of $2x$ is zero (it is $2 \times [\text{generator}]$ in a $\Zt$ module), but the $p = 3$ layer is nonzero:
\begin{equation}
2x = y \neq 0 \in E_\infty^{3,1}.
\end{equation}
This means $2 \times [p+ip] = [\text{CF}]$, and $4 \times [p+ip] = 2 \times [\text{CF}] = 0$.
The classification group is $\Zn{4}$, not $\Zt \times \Zt$.

\paragraph{Physical interpretation.}
Stacking two copies of the $p+ip$ superconductor decoration creates a Majorana chain decoration (through a cross-term), which then manifests as a nontrivial complex fermion class after purification.
This is a purely \emph{interaction-enabled} phenomenon: in the free-fermion classification, two copies of $p+ip$ would simply give a trivial state.

%% ============================================================
%% CHAPTER 7
%% ============================================================
\section{Computing Stacking Groups with SptSet}
\label{sec:sptset-stacking}

\subsection{From layers to the full group}
\label{subsec:layers-to-group}

The function that computes the full classification group (including extensions) is:
\begin{verbatim}
M := SptSetSpecSeqResult(SS, 3, [1,2,3]);
SptSetFpZModuleCanonicalForm(M);
Display(M);
\end{verbatim}

Let us parse the arguments:
\begin{itemize}
    \item \texttt{SS}: the spectral sequence object.
    \item \texttt{3}: the total degree $n = d + 1$ (here $d = 2$).
    \item \texttt{[1,2,3]}: the range of $p$ values to include ($p = 1, 2, 3$, corresponding to Majorana, CF, and bosonic layers).
\end{itemize}

The function iterates over the $p$ range and repeatedly applies the extension constructor \texttt{SptSetSpecSeqModuleExtension}, building up the full classification group from the graded pieces.

\subsection{The extension script for 2D wallpaper groups}
\label{subsec:ext-script}

The complete script for computing stacking groups is:
\begin{verbatim}
LoadPackage("SptSet");

for it in [2..17] do
  Display(it);
  SG := SpaceGroupBBNWZ(2, it);
  fSG := IsomorphismPcpGroup(SG);
  SG1 := Image(fSG);
  R := ResolutionAlmostCrystalGroup(SG1, 6);
  gs := GeneratorsOfGroup(SG);

  f := GroupHomomorphismByImagesNC(SG1, GL(1, Integers),
    List(gs, x -> Image(fSG, x)),
    List(gs, x -> [[DeterminantMat(x)]]));

  SS := FermionEZSPTSpecSeq(R, f);

  # Classification layers (E_infinity pages)
  layers := FermionSPTLayers(SS, 2);
  Display("Layers:");
  Display(layers);

  # Full stacking group (with extensions)
  M := SptSetSpecSeqResult(SS, 3, [1,2,3]);
  SptSetFpZModuleCanonicalForm(M);
  Display(M);
od;
\end{verbatim}

The output for each wallpaper group shows two things:
\begin{enumerate}
    \item \textbf{Layers}: the list of $E_\infty^{p,q}$ groups (graded pieces).
    \item \textbf{M}: the full classification group after solving the extension problem.
\end{enumerate}
Comparing these two reveals whether any nontrivial extensions occur.

\subsection{Reading the output: layers vs.\ full group}
\label{subsec:layers-vs-full}

Consider a wallpaper group where the output is:
\begin{verbatim}
Layers:
[ <ZL-Module with torsions [ 2 ]>,         # E_inf^{1,2} = Z_2
  <ZL-Module with torsions [ 2, 2 ]>,      # E_inf^{2,1} = Z_2^2
  <ZL-Module with torsions [ 2, 2, 4 ]> ]  # E_inf^{3,0} = Z_2^2 x Z_4

M:
<ZL-Module with torsions [ 2, 4, 4, 4 ]>   # full group = Z_2 x Z_4^3
\end{verbatim}

If the extension were trivial, the full group would be $\Zt \times \Zt^2 \times (\Zt^2 \times \Zn{4}) = \Zt^5 \times \Zn{4}$.
But the actual group is $\Zt \times \Zn{4}^3$, which is different.
The discrepancy reveals that nontrivial extensions occur between the layers, promoted by cross-terms in the stacking rules.

\subsection{Interpreting the results physically}
\label{subsec:interpret-results}

The full stacking group $\mathcal{C}$ answers the following physical questions:
\begin{enumerate}[label=(\roman*)]
    \item \textbf{How many distinct phases?} The order of $\mathcal{C}$ gives the total number.
    \item \textbf{What are the generators?} Each generator of $\mathcal{C}$ corresponds to a ``root phase''---a phase that cannot be obtained by stacking simpler ones.
    \item \textbf{What is the order of each generator?} This tells you how many copies must be stacked to obtain the trivial phase. For example, if a generator has order 4, stacking 4 copies gives the trivial state.
    \item \textbf{Are there unexpected relations?} A nontrivial extension means that stacking certain phases produces phases in a \emph{different} layer. For instance, stacking two Majorana phases might produce a complex fermion phase.
\end{enumerate}

\subsection{Known limitations}
\label{subsec:limitations}

The current implementation has several known limitations:
\begin{itemize}
    \item \textbf{Higher-page differentials}: The AHSS engine fully implements $d_2$ and $d_3$, but $d_r$ for $r \geq 4$ is not implemented. This is sufficient for 2+1D but may miss effects in 3+1D and beyond.
    \item \textbf{Purification beyond $r = 2$}: The class purification algorithm (which removes coboundary-trivial ambiguities) currently works only at $r = 2$. For higher pages, it raises an error.
    \item \textbf{$p + ip$ stacking}: Some stacking twisters involving the $p + ip$ layer in 3+1D are set to zero (placeholders). The $p + ip$ differentials $d_3$ and $d_4$ are currently unknown.
    \item \textbf{Assertion failures}: For certain wallpaper groups, the extension computation encounters assertion failures related to non-integer Bockstein images, indicating subtle bugs in the stacking law or purification.
\end{itemize}


%% ============================================================
%% CHAPTER 8
%% ============================================================
\section{Summary and Outlook}
\label{sec:outlook}

\subsection{What we have learned}

This lecture presented a systematic framework for the classification and stacking of interacting fermionic SPT phases:

\begin{enumerate}
    \item \textbf{Domain wall decoration} provides a constructive, physical picture: every FSPT phase corresponds to a specific pattern of lower-dimensional topological decorations on symmetry domain walls.
    
    \item \textbf{The AHSS} organizes the classification into a computable algebraic framework:
    \begin{itemize}
        \item The $E_2$ page gives the ``first approximation'' from independent group cohomology.
        \item Differentials $d_r$ encode obstructions and equivalences between layers.
        \item The $E_\infty$ page gives the graded pieces of the classification.
    \end{itemize}
    
    \item \textbf{The extension problem} captures the nontrivial stacking structure: cross-terms between layers mean that the classification group is generally not a direct product.
    The add-twisters encode the extension cocycles.
    
    \item \textbf{SptSet} implements the entire pipeline in GAP, making systematic classification possible for any symmetry group with a computable resolution.
\end{enumerate}

\subsection{Comparison with free-fermion classification}

For crystalline symmetries, the free-fermion classification comes from $K$-theory applied to the band structure.
The interacting classification from AHSS can differ in several ways:
\begin{itemize}
    \item \textbf{Collapse}: A $\Z$ free-fermion invariant may become $\Z_n$ when interactions are included (fewer distinct phases).
    \item \textbf{New phases}: The AHSS may find phases (especially at higher $q$) that have no free-fermion counterpart.
    \item \textbf{Modified group structure}: Even when the total number of phases is the same, the stacking group structure (extensions vs.\ direct products) can differ.
\end{itemize}

The domain wall decoration / AHSS approach naturally includes \emph{all} interacting effects from the start, rather than treating interactions as a perturbation on top of the free-fermion classification.

\subsection{Open problems and future directions}

\begin{enumerate}
    \item \textbf{Full 3+1D classification}: Completing the higher-page differentials ($d_r$ for $r \geq 4$) and the missing stacking twisters for the $p + ip$ layer.
    
    \item \textbf{4+1D and beyond}: Extending the spectrum and the engine infrastructure to handle higher-dimensional cases, where new types of decorations appear.
    
    \item \textbf{Anomalous SET phases}: Extending the framework to symmetry-enriched topological phases, where the bulk carries topological order.
    
    \item \textbf{Fermionic crystalline equivalence principle}: Using the equivalence principle to relate internal FSPT to crystalline FSPT, extending the classification to all 230 space groups systematically.
    
    \item \textbf{Physical observables}: Connecting the mathematical invariants to measurable quantities (surface states, response functions, entanglement spectra).
\end{enumerate}


%% ============================================================
%% APPENDICES
%% ============================================================
\appendix

\section{Group Cohomology Conventions}
\label{ap:group-cohomology}

\subsection{Inhomogeneous cochains}

An \emph{inhomogeneous $p$-cochain} on a group $G$ with values in a $G$-module $A$ is a function
\begin{equation}
f: \underbrace{G \times \cdots \times G}_{p} \longrightarrow A.
\end{equation}
The \emph{inhomogeneous coboundary} $\dd f \in C^{p+1}(G, A)$ is
\begin{align}
(\dd f)(g_1, \ldots, g_{p+1}) &= g_1 \cdot f(g_2, \ldots, g_{p+1}) \nonumber \\
&\quad + \sum_{i=1}^{p} (-1)^i f(g_1, \ldots, g_i g_{i+1}, \ldots, g_{p+1}) \nonumber \\
&\quad + (-1)^{p+1} f(g_1, \ldots, g_p). \label{eq:inhomo-coboundary}
\end{align}
One verifies $\dd \circ \dd = 0$.

\subsection{Cup product and higher cup products}

The \emph{cup product} of an inhomogeneous $p$-cochain $f$ and a $q$-cochain $g$ is the $(p+q)$-cochain
\begin{equation}
(f \cup g)(g_1, \ldots, g_{p+q}) := f(g_1, \ldots, g_p) \cdot g(g_{p+1}, \ldots, g_{p+q}).
\end{equation}
(This is Steenrod's $\cup_0$ product in the homogeneous notation.)

The \emph{higher cup products} $\cup_i$ for $i \geq 1$ are defined on homogeneous cochains and satisfy the relation
\begin{equation}
\dd(f \cup_i g) = (\dd f) \cup_i g + (-1)^{|f|} f \cup_i (\dd g) + (-1)^{|f|+|g|+i} (f \cup_{i-1} g + (-1)^{(|f|+i)(|g|+i)} g \cup_{i-1} f).
\end{equation}
These higher cup products are essential for writing the FSPT coboundary equations and stacking rules.
In SptSet, they are implemented as \texttt{Cup0@}, \texttt{Cup1@}, \texttt{Cup2@}, etc.


\section{GAP Quick Reference}
\label{ap:gap-quick}

\subsection{Essentials}
\begin{verbatim}
x := 5;                         # variable assignment (use := not =)
[1, 2, 3];                      # a list (1-indexed!)
rec(a := 1, b := "hello");      # a record
f := x -> x^2;                  # lambda function
\end{verbatim}

\subsection{Groups}
\begin{verbatim}
G := CyclicGroup(4);             # Z_4
G := DihedralGroup(8);           # D_4
G := SymmetricGroup(4);          # S_4
GeneratorsOfGroup(G);            # list of generators
Order(G);                        # group order
\end{verbatim}

\subsection{SptSet-specific}
\begin{verbatim}
LoadPackage("SptSet");           # load the package
SS := FermionEZSPTSpecSeq(R,f);  # construct spectral sequence
FermionSPTLayersVerbose(SS, d);  # print E_r pages
M := SptSetSpecSeqResult(SS, n, pRange);  # full stacking group
SptSetFpZModuleCanonicalForm(M); # canonical form
\end{verbatim}


\bibliography{lecture_note}

\end{document}
